(1) 因 \(\nabla\cdot\boldsymbol v=\nabla\cdot\boldsymbol B=0\)，感应方程化为 \(D\boldsymbol B/Dt=(\boldsymbol B\cdot\nabla)\boldsymbol v\)。均匀初场始终均匀，故 \(\dot B_x=SB_y,\ \dot B_y=\dot B_z=0\)，
\[
\boldsymbol B(t)=(B_{x0}+StB_{y0},B_{y0},B_{z0}).
\]
理想近似要求特征磁雷诺数 \(Rm=UL/\eta_m\gg1\)，\(\eta_m=1/(\mu_0\sigma)\)，且外部剪切可维持。
(2) 质点满足 \(x=X+StY,\ y=Y,\ z=Z\)，所以
\[
F=\begin{pmatrix}1&St&0\\0&1&0\\0&0&1\end{pmatrix},\qquad \det F=1.
\]
冻结定理对任意随流面积元给 \(\boldsymbol B=F\boldsymbol B_0/\det F\)，恰好重得第一问；\(\det F=1\) 同时保证不可压缩。
(3) 直接微分
\[
\dot u_B=\frac{\boldsymbol B\cdot\dot{\boldsymbol B}}{\mu_0}=\frac{S B_xB_y}{\mu_0}=S\sigma_{xy}.
\]
理想MHD磁应力 \(\sigma_{ij}=B_iB_j/\mu_0-\delta_{ij}B^2/(2\mu_0)\) 与速度梯度收缩为 \(\sigma_{ij}\partial_jv_i=S B_xB_y/\mu_0\)。均匀场 \(\nabla\times B=0\)，体内无Lorentz力；维持剪切的边界外力通过磁张力传功，不能误判为无源增能。
(4) \(B^2=(B_{x0}+StB_{y0})^2+B_{y0}^2+B_{z0}^2\)。当 \(S B_{y0}\ne0\)，极小在 \(t_*=-B_{x0}/(SB_{y0})\)，值 \(u_{B,\min}=(B_{y0}^2+B_{z0}^2)/(2\mu_0)\)；若只考察 \(t\ge0\) 且 \(t_*<0\)，最小在初始时刻。\(S=0\) 或 \(B_{y0}=0\) 则场不放大。长期 \(B_x\propto t\) 导致磁张力反作用增长，固定剪切外驱与理想均匀场近似最终须重新检验。
